% Generated draft: requires mathematical and source review.
\documentclass[11pt]{article}
\usepackage[margin=25mm]{geometry}
\usepackage{fontspec}
\setmainfont{Latin Modern Roman}
\usepackage{amsmath,amssymb,mathtools,unicode-math}
\setmathfont{Latin Modern Math}
\usepackage{enumitem,needspace,xurl,hyperref}
\setlength{\parindent}{0pt}
\setlength{\parskip}{5pt}
\setlength{\emergencystretch}{3em}
\newcommand{\N}{\mathbb{N}}
\newcommand{\Z}{\mathbb{Z}}
\newcommand{\Q}{\mathbb{Q}}
\newcommand{\R}{\mathbb{R}}
\newcommand{\C}{\mathbb{C}}
\newcommand{\F}{\mathbb{F}}
\newcommand{\A}{\mathbb{A}}
\newcommand{\PP}{\mathbb{P}}
\newcommand{\E}{\mathbb{E}}
\newcommand{\eps}{\varepsilon}
\newcommand{\dd}{\,\mathrm{d}}
\begin{document}
\section*{UnsolvedMath ID: 6000012}
\textbf{Source problem code:} catalogue entry
\subsection{Definitions and mathematical statement}
Let $M$ be the Riemann surface in the source, regarded as a real surface, with its intended projectively flat connection. For a regular $C^2$ closed curve $\gamma:S^1\to M$ and a real torsion-free affine representative $\nabla$, a standard projective inflection condition is
\[
\dot\gamma(t)\wedge\nabla_{\dot\gamma(t)}\dot\gamma(t)=0.
\]
Equivalently, the acceleration is tangent to the curve. This condition is unchanged by replacing $\nabla$ with a projectively equivalent torsion-free connection, since that replacement adds a tangential term to the acceleration. In a local real projective chart it is the ordinary inflection condition relative to straight geodesics.

Estimate the minimum possible number of inflection points of a closed curve in the stated geometry. For a specified admissible curve class $\mathcal C$, this would ask for a sharp lower bound on
\[
\inf_{\gamma\in\mathcal C}|\{t\in S^1:\gamma\text{ has an inflection at }t\}|,
\]
with an appropriate multiplicity convention if required. The source does not specify $\mathcal C$, embeddedness, homotopy restrictions, genericity or isolated inflections. It also does not distinguish a real projective connection from a holomorphic projective structure. Thus the displayed real-connection definition explains one conventional interpretation, but the record does not supply a unique universal numerical bound or all hypotheses needed for one.
\subsection{Short English statement}
Find a sharp minimum number of inflections for closed curves on a surface with projectively flat geometry. The connection type, allowed curves, and counting convention are unspecified.
\subsection{Sources}
\begin{itemize}
\item Furuhata, Matsuzoe, and Urakawa, \emph{Open Problems in Affine Differential Geometry and Related Topics} (1998), item 4(a), scanned PDF pages 1--2, \url{https://www.jstage.jst.go.jp/article/iis/4/2/4_2_125/_pdf/-char/en}; accessed 2026-07-29, manual OCR from page images recorded in the provenance. The record attributes the question to Takeshi Sasaki.
\item The supplied summary mentions flex theorems and Pl\"ucker-type identities as partial context, without specifying publications or their hypotheses.
\item Local registry \texttt{data/unsolvedmath/problems.json}, permanent UnsolvedMath ID $6000012$, source code catalogue entry, \emph{Inflection Points of Projectively Flat Curves}. Dataset attribution: Ulam AI and the UnsolvedMath Contributors.
\item The supplied rights note says the J-STAGE article page marks this article CC BY 4.0.
\end{itemize}
\end{document}
